% docs/translation/latex/chapters/user_differences.tex - chapter of the LaTeX edition of the book "translation_report" (SIMC++ documentation).
%   Converted once from docs/translation/markdown/user_differences.md by tools/md2latex.py; from then on THIS FILE is the source of the
%   LaTeX edition and may be edited by hand. Re-running tools/md2latex.py --write-sources for this
%   chapter OVERWRITES the hand edits (see docs/README.md).
%   Compile: cmake --build <build> --target pdf, or tools/build_docs_pdf.py (figures, xelatex passes,
%   makeindex with docs/common/latex/simcpp.ist); by hand, see docs/README.md. XeLaTeX, class
%   docs/common/latex/simcpp.cls (packages: fontspec, tikz, xcolor, hyperref, listings, longtable, booktabs,
%   titlesec, titletoc, fancyhdr, makeidx, preview, ...); fonts TeX Gyre Pagella, TeX Gyre Heros, DejaVu Sans
%   Mono. The figures are built first from docs/{manual,translation}/figures/*.tex (preview.sty) into PDFs.
\chapter{\texorpdfstring{Differences between SIMC++ and the Fortran SIMC (user guide)}{Differences between SIMC++ and the Fortran SIMC (user guide)}}\label{translation_report-useru5fdifferences}
\simcabstract{The complete, line-by-line list is \code{docs/\allowbreak{}translation/\allowbreak{}markdown/\allowbreak{}FORTRAN\_\allowbreak{}DIFFERENCES.\allowbreak{}md}. Each item there gives the Fortran file and line, the
original behaviour, the C++ behaviour, the measured effect and the test that covers it. This guide summarises it for
users: \textbf{what changes in your numbers, and when}. Everything in 10.1 is undone by \code{-{}-{}fortran-{}compat}\simcidx{Command-line options!fortran-compat@\code{-{}-{}fortran-{}compat}}.}
\section*{\texorpdfstring{1. Corrections of the default mode (Category A)}{1. Corrections of the default mode (Category A)}}\addcontentsline{toc}{section}{\texorpdfstring{1. Corrections of the default mode (Category A)}{1. Corrections of the default mode (Category A)}}\label{translation_report-useru5fdifferences:1-corrections-of-the-default-mode-category-a}
\subsection*{\texorpdfstring{Changes that can move yields, weights or acceptances}{Changes that can move yields, weights or acceptances}}\label{translation_report-useru5fdifferences:changes-that-can-move-yields-weights-or-acceptances}
\begin{simctable}{footnotesize}{3}{>{\raggedright\arraybackslash}p{\dimexpr 0.190\linewidth-2\tabcolsep\relax}>{\raggedright\arraybackslash}p{\dimexpr 0.299\linewidth-2\tabcolsep\relax}>{\raggedright\arraybackslash}p{\dimexpr 0.449\linewidth-2\tabcolsep\relax}}
\caption{Changes that can move yields, weights or acceptances}\\%
\toprule\rowcolor{simchead}\simcth{id} & \simcth{when it matters} & \simcth{what changes in the default mode} \\\midrule\endfirsthead
\multicolumn{3}{@{}l}{\simccontinued}\\\toprule\rowcolor{simchead}\simcth{id} & \simcth{when it matters} & \simcth{what changes in the default mode} \\\midrule\endhead
\midrule\multicolumn{3}{r@{}}{\simccontinues}\\\endfoot
\bottomrule\endlastfoot
D-006 pass\_cuts & every run where the reconstructed hadron delta can lie between \code{SPedge\%\allowbreak{}p\%\allowbreak{}delta\%\allowbreak{}max}\simcidx{Parameters!SPedge.p.delta.max@\code{SPedge\%p\%delta\%max}} and \code{SPedge\%\allowbreak{}e\%\allowbreak{}delta\%\allowbreak{}max}\simcidx{Parameters!SPedge.e.delta.max@\code{SPedge\%e\%delta\%max}} & the Fortran tests the hadron delta against the \textbf{electron} delta maximum. The default mode uses the hadron limit. This changes \code{wtcontr}, the resolutions,\allowbreak{} and with \code{hard\_\allowbreak{}cuts}\simcidx{Parameters!hard_cuts@\code{hard\_cuts}} the accepted events: −0.01 \% to −0.23 \% in the K3 H(e,e'p) runs, −24 \% to −34 \% of ΣWeight in rho runs, up to −24 \% in spect\_mode runs, 0 when the hadron window lies inside the electron one \\
D-043 F5, D-051 F6 survival momentum & kaon and SIDIS runs with \code{doing\_\allowbreak{}decay\ =\ 0} & the survival probability uses the hadron's own momentum, not the central momentum: the kaon \textbf{weights} (and K± SIDIS weights) and the \code{decdist} column \\
D-043 F6, D-051 F6 aerogel distance & decay-OFF kaon/SIDIS runs with an HMS hadron (path length with the missing square root); SHMS unaffected & survival probability \\
D-061 F1, F1e HRS & every HRS-L/HRS-R run & the intended transport lengths (Fortran: path length −3e232 cm, "decays" everywhere, kaon survival weight +Inf) and the 45° VDC rotation (Fortran: no rotation). HRS acceptance, FP path, decay \\
D-048 F2 Delta tail 1 & radiative Delta runs & the incoming electron radiates (the Fortran fails every such try): Ntried, wtcontr \\
D-048 F1 Delta D/A & Delta runs on D or A & they run (the Fortran always STOPs); there is no Fortran reference for them \\
D-033 F3 Benhar SF & A(e,e'p) with \code{use\_\allowbreak{}benhar\_\allowbreak{}sf}\simcidx{Parameters!use_benhar_sf@\code{use\_benhar\_sf}} & the proper differential spectral function of each node instead of a bin probability divided by the event's own Pm² (divergent at small Pm): weights, wtcontr \\
D-051 F1, F2 SIDIS & SIDIS events with all PDFs 0, or D/A runs with \code{do\_\allowbreak{}fermi}\simcidx{Parameters!do_fermi@\code{do\_fermi}} and a negative x\_Bj & weight 0 instead of NaN; the run no longer STOPs \\
D-051 F4 π0 SIDIS & π0 SIDIS & the d̄ fragmentation function and the π0 mass \\
D-057 F1, F6, F9, F21 rho & rho runs & energy loss of the ρ with its own angles; efer of D/³He; yptar sign of far-side pions; a NaN pion (1 event per 1e5) \\
D-057 F2 rho decay & rho runs & no decay after a failed try: no weight change, but the \textbf{random stream} differs from the Fortran's \\
D-037 F4, F7 collimator / decay & pion/kaon runs with collimators or decay & a particle that ranges out in the collimator stops; the track after a decay ends at the right z (≤ 0.28 cm in the path length, no acceptance decision changed in the captures) \\
D-062 G1, G2, G3 π0 / calorimeter & π0 runs with \code{doing\_\allowbreak{}decay\ =\ 1}; calorimeter as electron arm with decay; π0 photon path & photons (and an electron) never decay, so the run works (the Fortran STOPs); each photon path starts at 0 \\
D-068 π0 in HMS/SHMS with decay & π0 transported in the HMS/SHMS with \code{doing\_\allowbreak{}decay\ =\ 1} & the π0 is stable (the Fortran STOPs) \\
D-067 T1, T2 target field & \code{using\_\allowbreak{}tgt\_\allowbreak{}field}\simcidx{Parameters!using_tgt_field@\code{using\_tgt\_field}} with \code{targ\_\allowbreak{}Bphi}\simcidx{Parameters!targ_Bphi@\code{targ\_Bphi}} not on an arm's side, or field on the hadron side with targ\_\allowbreak{}Bangle < θp & the field angles of the arms \\
D-033 F7, D-028 F2 NaN guards & rare edge events (A(e,e'p) extreme edges; H elastic rounding) & NaN events are rejected or clamped instead of carrying NaN values \\
D-015 F9 raddeg & HMS/SHMS rotated apertures & π/180 exactly instead of 8 digits (no decision changed in the validation runs) \\
D-017 extreme trip & pudding can or cryo17 with a hadron angle > 45° & the hadron energy-loss limits (e.g. 1.3565 → 1.3591 MeV) \\
\end{simctable}
\subsection*{\texorpdfstring{Changes in tree columns only (no effect on weights)}{Changes in tree columns only (no effect on weights)}}\label{translation_report-useru5fdifferences:changes-in-tree-columns-only-no-effect-on-weights}
\begin{simctable}{small}{3}{>{\raggedright\arraybackslash}p{\dimexpr 0.167\linewidth-2\tabcolsep\relax}>{\raggedright\arraybackslash}p{\dimexpr 0.287\linewidth-2\tabcolsep\relax}>{\raggedright\arraybackslash}p{\dimexpr 0.478\linewidth-2\tabcolsep\relax}}
%
\toprule\rowcolor{simchead}\simcth{id} & \simcth{column} & \simcth{change} \\\midrule\endfirsthead
\multicolumn{3}{@{}l}{\simccontinued}\\\toprule\rowcolor{simchead}\simcth{id} & \simcth{column} & \simcth{change} \\\midrule\endhead
\midrule\multicolumn{3}{r@{}}{\simccontinues}\\\endfoot
\bottomrule\endlastfoot
D-051 F3 & SIDIS/rho \code{pfermi} & 0 instead of NaN for H and A ≥ 3 \\
D-043 F1, F2, F3 & kaon \code{saghai} & the Saghai cross section (the Fortran writes 0 in every event; its tables were also read off by one line and on a wrong s grid) \\
D-048 F3 & Delta \code{sigcm} & a meaningful pion fit value (the Fortran's is garbage, NaN in 0.4 \% of the events) \\
D-067 T14 & π0 and rho runs with the target field & each value in its own column (the Fortran overwrites the polarized-angle columns) \\
D-006 kaon decay & \code{Mhadron} after a K → ππ decay in the first half of a magnet & the pion mass instead of the kaon mass \\
D-038 F10, D-033 F5, D-051 F9, D-057 F3 & internal values (\code{sigcm2} above W = 2 GeV, vertex PmPer, Mh2\_final) & not written to any output \\
\end{simctable}
\subsection*{\texorpdfstring{Changes in the text reports only}{Changes in the text reports only}}\label{translation_report-useru5fdifferences:changes-in-the-text-reports-only}
\begin{simctable}{footnotesize}{3}{>{\raggedright\arraybackslash}p{\dimexpr 0.182\linewidth-2\tabcolsep\relax}>{\raggedright\arraybackslash}p{\dimexpr 0.221\linewidth-2\tabcolsep\relax}>{\raggedright\arraybackslash}p{\dimexpr 0.535\linewidth-2\tabcolsep\relax}}
%
\toprule\rowcolor{simchead}\simcth{id} & \simcth{file} & \simcth{change} \\\midrule\endfirsthead
\multicolumn{3}{@{}l}{\simccontinued}\\\toprule\rowcolor{simchead}\simcth{id} & \simcth{file} & \simcth{change} \\\midrule\endhead
\midrule\multicolumn{3}{r@{}}{\simccontinues}\\\endfoot
\bottomrule\endlastfoot
D-038 F5, D-043 F13, D-048 F9 & \code{.\allowbreak{}hist} \code{CENTRAL.\allowbreak{}sigcc} of pion, kaon and Delta runs & computed at the real central kinematics (the Fortran's is ×4.9 to ×11 off). A report value; not used anywhere else \\
D-011, D-060, D-061 F8, D-037 F2 & \code{.\allowbreak{}geni} & SHMS "Q3 mid" shows q3\_mid (Fortran: q2\_mid); SOS/HRS slit hor/vert in label order; the SHMS "DC1, DC2" label; HRS "Successes" lines; a "Collimator stops (pi/mu)" line \\
D-021, D-037 F1 & \code{.\allowbreak{}geni} HMS/SHMS slit counters with collimators & each slit counter counted once per ray (Fortran: once per collimator step) \\
D-018 & \code{.\allowbreak{}hist} & integers written as integers (Fortran: T/F); angle offsets in mr; the Delta title \\
D-038 F15, D-051 F12, F14, D-057 F15, F18 & \code{.\allowbreak{}hist} titles; rho \code{.\allowbreak{}geni}/\code{.\allowbreak{}gen}/\code{.\allowbreak{}hist} P-arm lines & titles for coherent π−, π0 SIDIS, A ≥ 3 kaon SIDIS (the Fortran STOPs), D/³He rho; rho P-arm histograms filled with the decay pion \\
D-014, D-048 F6 & log & the de Forest cos φ warning is printed; the Delta warning text \\
\end{simctable}
\subsection*{\texorpdfstring{Changes in what is accepted as input}{Changes in what is accepted as input}}\label{translation_report-useru5fdifferences:changes-in-what-is-accepted-as-input}
\begin{simctable}{small}{3}{>{\raggedright\arraybackslash}p{\dimexpr 0.147\linewidth-2\tabcolsep\relax}>{\raggedright\arraybackslash}p{\dimexpr 0.330\linewidth-2\tabcolsep\relax}>{\raggedright\arraybackslash}p{\dimexpr 0.456\linewidth-2\tabcolsep\relax}}
%
\toprule\rowcolor{simchead}\simcth{id} & \simcth{input} & \simcth{default mode} \\\midrule\endfirsthead
\multicolumn{3}{@{}l}{\simccontinued}\\\toprule\rowcolor{simchead}\simcth{id} & \simcth{input} & \simcth{default mode} \\\midrule\endhead
\midrule\multicolumn{3}{r@{}}{\simccontinues}\\\endfoot
\bottomrule\endlastfoot
D-032 & arm code 6 & rejected (both modes) \\
D-038 F1 & H(e,\allowbreak{}e'π−)Δ++ (\code{which\_\allowbreak{}pion\ =\ 3}) & runs (the Fortran STOPs) \\
D-038 F2 & π0 from the neutron on hydrogen & refused (the Fortran runs it with a neutron mass on a proton target) \\
D-037 F9 & \code{doing\_\allowbreak{}decay}\simcidx{Parameters!doing_decay@\code{doing\_decay}} with a hadron that is not π/K & refused (the Fortran STOPs silently at the first decay) \\
D-025/D-032 & SHMS LSA tune & checked at every call (no valid input affected) \\
\end{simctable}
\section*{\texorpdfstring{2. Kept as in the Fortran: acceptance mode (Category C)}{2. Kept as in the Fortran: acceptance mode (Category C)}}\addcontentsline{toc}{section}{\texorpdfstring{2. Kept as in the Fortran: acceptance mode (Category C)}{2. Kept as in the Fortran: acceptance mode (Category C)}}\label{translation_report-useru5fdifferences:2-kept-as-in-the-fortran-acceptance-mode-category-c}
The user has not decided this yet (D-031 item 1), so it is kept in both modes. On hydrogen, \code{doing\_\allowbreak{}phsp}\simcidx{Parameters!doing_phsp@\code{doing\_phsp}} keeps the
elastic kinematics; radiation stays on; the target mass for the photon limits is 0
(5.8\footnote{SIMC++ User Manual, ``5.8 Phase-space and acceptance runs (\code{doing\_\allowbreak{}phsp})'' (\code{docs/\allowbreak{}manual/\allowbreak{}markdown/\allowbreak{}user/\allowbreak{}05\_\allowbreak{}reactions.\allowbreak{}md})}).

\section*{\texorpdfstring{3. Questions open for experts (the expert list)}{3. Questions open for experts (the expert list)}}\addcontentsline{toc}{section}{\texorpdfstring{3. Questions open for experts (the expert list)}{3. Questions open for experts (the expert list)}}\label{translation_report-useru5fdifferences:3-questions-open-for-experts-the-expert-list}
These Fortran behaviours \textbf{might} be intentional. SIMC++ keeps them in both modes (\code{docs/\allowbreak{}translation/\allowbreak{}markdown/\allowbreak{}FORTRAN\_\allowbreak{}DIFFERENCES.\allowbreak{}md} §4).
In plain words:

\begin{itemize}
\item \textbf{Delta weight without cross section} (D-048 F4): Delta yields are not rates. A related question is whether a factor
2 is missing when both proton-energy roots are physical (F5).
\item \textbf{Radiative corrections} (D-019): a disabled tail can still be drawn with \code{one\_\allowbreak{}tail}\simcidx{Parameters!one_tail@\code{one\_tail}}; switching off the hadron tail
also removes the e-p interference term; with \code{intcor\_\allowbreak{}mode\ =\ 0} part of a term is lost; the input \code{lambda}\simcidx{Parameters!lambda@\code{lambda}} is
overwritten (D-005); a few stale values for specific flag settings.
\item \textbf{D(e,e'p)} has no hadron-angle jacobian (≤ 1.5 \%), A(e,e'p) does (D-033 F4). He-3/He-4 use the C-12 spectral
function (F1).
\item \textbf{Spectrometer details}: HMS hodoscope 2 tested with hodoscope-1 limits, a Cherenkov foil counted twice, SHMS drift
lengths that differ from the COSY values (D-015); the beer-can radius differs between two routines (D-017); the SOS last
7.62 cm has no path length or decay (D-060); SOS/HRS reconstruct with the raster position, HMS/SHMS without (D-061 F3);
the HRS aerogel distance (F6).
\item \textbf{Collimators and decay}: the collimator scattering ignores \code{mc\_\allowbreak{}smear}\simcidx{Parameters!mc_smear@\code{mc\_smear}}; the pion absorption parameters; where a decay
inside a magnet is applied (D-037 F3, F5, F8).
\item \textbf{Pion and kaon models}: coherent production uses nucleon models; the input final-nucleus mass is ignored for coherent
production; a new Δ mass when tail 1 recomputes the event; single-precision MAID amplitudes (D-038 F3, F4, F9, F12); a
fixed kaon decay model, the Saghai angles, the Saghai table records, the proton mass for Σ− (D-043).
\item \textbf{SIDIS}: the missing-mass threshold formula and its use for kaons and with Fermi motion, the generated hadron range
with radiation, R\_h for deuterium (D-051 F7, F8, F10, F16).
\item \textbf{Rho}: no Fermi motion for D/³He, s without Fermi terms, t at the pole mass, the decay frame, the xptar formula, the
radiation of the neutral ρ (70 \% of tail-3 tries fail) and a new ρ when tail 1 recomputes the event (D-057).
\item \textbf{Calorimeter and π0}: the calorimeter windows reuse the previous trip's values (D-012 F4); \code{drift\_\allowbreak{}to\_\allowbreak{}cal}\simcidx{Parameters!drift_to_cal@\code{drift\_to\_cal}} 0 is not
checked (D-062 G5); a second photon pair in tail 1 (G8); a π0 in the HMS/SHMS is transported as a charged particle (D-065
P1) and is bent in the target field (D-067 T13).
\item \textbf{Target field}: the step limit depends on the arm side, the convergence test, the 90/E constant, a polarization
vector without \code{targ\_\allowbreak{}Bphi}\simcidx{Parameters!targ_Bphi@\code{targ\_Bphi}} (D-067 T4-T6, T15).
\item \textbf{Bookkeeping}: with ngen < 0 wtcontr is divided by the tries; \code{doing\_\allowbreak{}herho} only for A = 3; SAVEd values in a few
report quantities (D-024, D-015 tmpfact, D-018).
\end{itemize}
\section*{\texorpdfstring{4. Deliberate differences of SIMC++ (Category DEV)}{4. Deliberate differences of SIMC++ (Category DEV)}}\addcontentsline{toc}{section}{\texorpdfstring{4. Deliberate differences of SIMC++ (Category DEV)}{4. Deliberate differences of SIMC++ (Category DEV)}}\label{translation_report-useru5fdifferences:4-deliberate-differences-of-simc-category-dev}
\begin{itemize}
\item Default generator mt64 instead of RANLUX; RANLUX available (8.2\footnote{SIMC++ User Manual, ``8.2 Random-number generators: \code{mt64} and \code{ranlux}'' (\code{docs/\allowbreak{}manual/\allowbreak{}markdown/\allowbreak{}user/\allowbreak{}08\_\allowbreak{}running\_\allowbreak{}and\_\allowbreak{}reproducibility.\allowbreak{}md})}); \code{random\_\allowbreak{}state\_\allowbreak{}file}\simcidx{Parameters!random_state_file@\code{random\_state\_file}} ignored.
\item STOP → exit status 3; rejected input → 2 (the Fortran always exits 0).
\item JSON input as a superset of the \code{.\allowbreak{}inp}; plug-in selection; \code{-{}-{}threads}\simcidx{Command-line options!threads@\code{-{}-{}threads}}.
\item Outputs: the ROOT tree written directly, plus the summary JSON and \code{.\allowbreak{}gen.\allowbreak{}root}; no \code{.\allowbreak{}bin}.
\item The model tables are read from the data directory at start-up and md5-checked by the tests; a missing file stops the run
before it begins.
\item The D-033 F1 spectral-function fallback warns and can be overridden from the JSON.
\item rho on A ≥ 4 and rho with another reaction flag or with \code{doing\_\allowbreak{}pizero}\simcidx{Parameters!doing_pizero@\code{doing\_pizero}} are refused at input.
\end{itemize}
